% Generated by blueprint/check.py. Do not edit.
\begin{longtable}{p{0.06\linewidth}p{0.16\linewidth}p{0.31\linewidth}p{0.35\linewidth}}
\textbf{ID} & \textbf{Class/source} & \textbf{Disposition} & \textbf{Public effect} \\
\hline
G01 & false-statement; \seqsplit{appendix-preliminaries.tex:34} \seqsplit{[13119]} & Prove that its coordinates sum to one, hence it is nonzero, and derive the probability estimate from this identity. & Proof repair only; no public QPBT hypothesis or conclusion changes. \\
G02 & index-omission; \seqsplit{appendix-strategies.tex:42-104,130} \seqsplit{[13237-13300,13325];} \seqsplit{appendix-combine-xz.tex:32} \seqsplit{[13750]} & Restore indicator indices, quantify the point domain, use X/Z bases, and enumerate all seven decider edges and both orientations. & Restores the intended tests without strengthening soundness. \\
G03 & missing-case; \seqsplit{appendix-strategies.tex:92-97} \seqsplit{[13287-13297]} & Use an internal small-degree helper and discharge the complementary public branch with a proved universal large-error bound. & The source-faithful A.7 and final soundness statements remain unconditional beyond admissibility. \\
G04 & typo-and-domain; \seqsplit{qpbt-game-and-soundness.tex:533-591} \seqsplit{[5579-5635]} & Use Alice/Bob suffixes consistently, one argument order, and the faithful boundary 0 <= epsilon <= 1. & Typographical normalization plus an explicit probability-domain boundary. \\
G05 & norm-mismatch; \seqsplit{appendix-separate-xz-conclude.tex:256-369,450-462} \seqsplit{[14709-14915]} & Prove a named square-root robustness bridge, positivity/contraction bounds, and small/large error cases; rechoose universal constants. & Only universal polynomial constants and exponent are reparameterized. \\
G06 & label-typo; \seqsplit{appendix-strategies.tex:112} \seqsplit{[13307]} & Preserve the literal source label in metadata and use a normalized Lean name. & Metadata only. \\
G07 & missing-lemma; \seqsplit{appendix-strategies.tex:8-16} \seqsplit{[13203-13211]} & Prove one simultaneous local Naimark dilation preserving every game correlation and compose it with extraction isometries. & Discharges an implicit proof boundary without adding an assumption. \\
G08 & missing-conditioning; \seqsplit{appendix-combine-xz.tex:117-153} \seqsplit{[13835-13871]} & Prove a good/bad fiber lemma with threshold eta and average error O(delta/eta + eta). & Internal robustness exponent loss only. \\
G09 & false-statement; \seqsplit{appendix-separate-xz-conclude.tex:37-43} \seqsplit{[14490-14496]} & State the exact trace phase and specialize only after proving that the phase vanishes. & A false intermediate is weakened; every downstream same-basis use is retained. \\
G10 & missing-hypotheses; \seqsplit{classical-ldt.tex:399-443} \seqsplit{[4561-4605]} & Expose admissibility and sampler dimension assumptions, correct the code tuple, and prove tensor-code/simultaneous-test adapters. & Narrows the literal internal theorem to the game actually defined. \\
G11 & dimension-mismatch; \seqsplit{appendix-apply-classical-ldt.tex:300-321} \seqsplit{[14306-14327]} & Prove a direct Fin n axis LDT and equivalence to the CL-indexed version when divisibility holds. & Preserves the QPBT admissibility assumptions. \\
G12 & malformed-references; \seqsplit{appendix-apply-classical-ldt.tex:71-275} \seqsplit{[14077-14281]} & Formalize all marginals and mixture weights, add Property 3, repair indices and expectations, and state the complement estimate. & Proof repair only. \\
G13 & missing-proof; \seqsplit{appendix-apply-classical-ldt.tex:432} \seqsplit{[14438]} & Prove a named Cauchy-Schwarz/projectivity bound and handle alpha = 0 separately. & Fills proof debt without changing A.21. \\
G14 & index-typos; \seqsplit{pauli.tex:123-148,191-196} \seqsplit{[2071-2096,2139-2144]} & Retain arbitrary natural L, range j through the k basis coordinates, and index the resulting single-qubit projector by a\_j. & Repairs two indices in the qudit-to-qubit conversion without changing its hypotheses or conclusion. \\
G15 & index-typo; \seqsplit{appendix-strategies.tex:198-203} \seqsplit{[13393-13398]} & Use Z\textasciicircum{}(r\_Z)(u\_Z) X\textasciicircum{}(r\_X)(u\_X) in the reversed product throughout the anticommutation argument. & Repairs an internal index typo; the stated observable consequence is unchanged. \\
\end{longtable}
